% ============================================================================
% Fichier  : partie6/P06_C20_mesurer_gouvernance_recherche.tex
% Titre    : Mesurer la Gouvernance des SI : État de la Recherche
% Auteur   : MINKA MI NGUIDJOÏ Thierry Emmanuel
% Version  : 1.0 -- octobre 2026
% ============================================================================
% Chapitre de recherche. Les références citées sont des publications
% académiques identifiées ; leurs résultats sont résumés, non reproduits.
% Les chiffres repris d'une étude sont attribués à leurs auteurs.
% ============================================================================

\chapter{Mesurer la Gouvernance des SI : État de la Recherche}
\label{chap:recherche-mesure}
\chaptermark{État de la recherche}

\epigraph{%
    « Ce qui compte ne peut pas toujours être compté,\\
    et ce qui peut être compté ne compte pas toujours. »%
}{Maxime attribuée, sous diverses formes, à plusieurs auteurs}

\niveauChapitre{Recherche}{%
    Maîtrise du Chapitre~\ref{chap:cobit-capacite} (capacité et maturité)
    et du Chapitre~\ref{chap:alignement} (alignement et valeur).%
}

\begin{boxobjectifs}
\begin{itemize}
    \item Synthétiser les travaux de référence sur la gouvernance des SI.
    \item Discuter la validité des modèles de maturité.
    \item Lire de façon critique une étude empirique.
    \item Formuler une question de recherche sur la gouvernance des SI en
          Afrique centrale.
\end{itemize}
\end{boxobjectifs}

\bigskip

% ============================================================================
\section*{Question directrice : la gouvernance des SI sert-elle à quelque chose~?}
% ============================================================================

Tout ce cours repose sur une conviction~: une meilleure gouvernance des SI
produit de meilleurs résultats. Le praticien l'admet volontiers~; le
chercheur demande des preuves. Qu'entend-on par « meilleure
gouvernance »~? Comment la mesure-t-on~? Et si l'on observe que les
organisations bien gouvernées réussissent mieux, est-ce la gouvernance qui
produit le succès, ou le succès qui permet de financer la gouvernance~?

La Partie~\ref{part:recherche}, qui s'ouvre ici, prend ce recul. Ce
chapitre présente les grandes étapes de la recherche
(section~\ref{sec:c20-courants}), ce qu'elle a établi
(section~\ref{sec:c20-resultats}), le débat sur les modèles de maturité
(section~\ref{sec:c20-maturite}), une méthode de lecture critique
(section~\ref{sec:c20-lecture}), puis les pistes de recherche ouvertes en
Afrique centrale (section~\ref{sec:c20-afrique}). Il s'adresse en
particulier aux étudiants qui préparent un mémoire.

\begin{boxinfo}
\textbf{Règle de lecture.} Les travaux cités dans ce chapitre sont des
publications académiques identifiées dans la bibliographie. Leurs
résultats sont résumés et attribués à leurs auteurs. Un résumé n'est pas
une lecture~: avant de citer une étude dans un mémoire, lisez-la.
\end{boxinfo}

\separateur

% ============================================================================
\section{Trois Courants de Recherche}
\label{sec:c20-courants}
% ============================================================================

\subsection{Les structures de décision}

Le premier courant s'est demandé \emph{où} se prennent les décisions
informatiques. Dans les années 1990, la recherche oppose les modèles
centralisés, décentralisés et fédéraux. Sambamurthy et Zmud montrent que
le bon arrangement dépend de plusieurs facteurs de contingence qui
agissent ensemble~: la taille, la stratégie, la structure de
l'organisation~\cite{sambamurthy1999arrangements}. Weill et Ross
prolongent ce travail avec les archétypes et la matrice des droits de
décision (Chapitre~\ref{chap:alignement})~\cite{weill2004itgovernance}.

\subsection{Les mécanismes}

Le deuxième courant passe de la question « qui décide~? » à la question
« par quels moyens~? ». Peterson distingue les structures, les processus
et les capacités relationnelles~\cite{peterson2004crafting}~; De Haes et
Van Grembergen reprennent et outillent cette distinction en structures,
processus et mécanismes relationnels (Chapitre~\ref{chap:gouv-entreprise}),
et étudient leur effet sur l'alignement~\cite{dehaes2009exploratory,dehaes2015egit}.

\subsection{L'effet sur la performance}

Le troisième courant cherche à établir le lien entre gouvernance des SI,
alignement et performance de l'organisation. Brown et Grant en ont dressé
une revue structurante au milieu des années 2000~\cite{brown2005framing}.
Les études quantitatives se sont depuis multipliées, avec des méthodes de
plus en plus exigeantes, comme celle de Wu, Straub et Liang, qui croise
les réponses de dirigeants métiers et de responsables informatiques des
mêmes entreprises~\cite{wu2015mechanisms}.

\begin{tableaucours}{L{2.6cm}L{4.8cm}L{4.8cm}}{%
    \tableauHeader{Courant} &
    \tableauHeader{Question} &
    \tableauHeader{Travaux de référence}
}
Structures & Où se prennent les décisions~? & Sambamurthy et Zmud (1999)~;
Weill et Ross (2004) \\
Mécanismes & Par quels moyens la gouvernance s'exerce-t-elle~? & Peterson
(2004)~; De Haes et Van Grembergen (2009, 2015) \\
Performance & La gouvernance améliore-t-elle les résultats~? & Brown et
Grant (2005)~; Wu, Straub et Liang (2015) \\
\end{tableaucours}
\vspace{-0.8em}
\begin{center}{\small\itshape Tableau~20.1. Trois courants de la recherche
sur la gouvernance des SI}\end{center}

% ============================================================================
\section{Ce que la Recherche a Établi}
\label{sec:c20-resultats}
% ============================================================================

\subsection{Des résultats convergents}

Avec les précautions qui s'imposent, trois résultats reviennent avec
régularité dans la littérature.

\begin{enumerate}
    \item \textbf{La gouvernance compte.} Weill et Ross ont rapporté, sur
          un large échantillon d'entreprises, que celles dont la
          gouvernance des SI était la plus efficace obtenaient, à
          stratégie comparable, une rentabilité nettement supérieure à
          celles dont la gouvernance était la plus faible, de l'ordre de
          20~\% selon leurs mesures~\cite{weill2004itgovernance}.
    \item \textbf{Les mécanismes relationnels comptent autant que les
          structures.} Plusieurs études observent qu'à structures égales,
          ce sont la compréhension mutuelle et la collaboration entre
          métiers et informatique qui font la différence
          (Chapitre~\ref{chap:cas-integrateurs}~: le cas de la BANQUE SANAGA un an
          après).
    \item \textbf{L'effet passe par l'alignement.} La gouvernance agit sur
          la performance surtout par l'intermédiaire de l'alignement
          stratégique~: elle aide à faire les bons choix plus qu'elle ne
          produit directement des résultats~\cite{wu2015mechanisms}.
\end{enumerate}

\subsection{Ce qui reste débattu}

\begin{tableaucours}{L{4.6cm}L{7.6cm}}{%
    \tableauHeader{Question} &
    \tableauHeader{État du débat}
}
Le sens de la causalité & Les organisations performantes ont-elles une
bonne gouvernance parce qu'elles en ont les moyens~? Les études
transversales ne permettent pas de trancher. \\
La mesure de la gouvernance & Elle repose souvent sur l'autoévaluation
d'un seul répondant, ce qui fragilise les résultats. \\
La transposition des résultats & La plupart des études portent sur des
organisations d'Amérique du Nord, d'Europe et d'Asie~; leur validité
ailleurs n'est pas démontrée. \\
L'effet des référentiels & L'adoption de COBIT ou d'ITIL est fréquente,
mais leur effet propre, toutes choses égales par ailleurs, reste mal
établi. \\
\end{tableaucours}
\vspace{-0.8em}
\begin{center}{\small\itshape Tableau~20.2. Questions débattues}\end{center}

\begin{boxalert}
\textbf{Violation typique.} Un mémoire affirme que « COBIT améliore la
performance des entreprises de 20~\% », en citant Weill et Ross. Or Weill
et Ross n'étudient pas COBIT, et leur écart de rentabilité est une
association mesurée sur leur échantillon, pas un effet causal
transposable.

\smallskip
\textbf{Règle~:} citer un résultat, c'est citer aussi ce qu'il mesure, sur
quel échantillon, et avec quelle méthode.
\end{boxalert}

% ============================================================================
\section{Le Débat sur les Modèles de Maturité}
\label{sec:c20-maturite}
% ============================================================================

\subsection{Un succès considérable}

Depuis le modèle de maturité du logiciel de l'université Carnegie Mellon,
les modèles de maturité se sont multipliés~: COBIT en utilise un depuis
sa version~3, et le modèle de capacité de COBIT~2019 s'inspire du CMMI
(Chapitre~\ref{chap:cobit-capacite}). Leur succès tient à leur
simplicité~: une échelle de 0 à~5, que tout dirigeant comprend.

\subsection{Quatre critiques}

La recherche en a pourtant relevé les limites~\cite{becker2009maturity,poppelbuss2011maturity}.

\begin{enumerate}
    \item \textbf{La linéarité.} Les modèles supposent un chemin unique,
          du niveau~1 au niveau~5. Or une organisation peut avoir de bonnes
          raisons de ne pas viser le niveau~5, ou de progresser autrement.
    \item \textbf{Le fondement théorique.} Beaucoup de modèles sont construits
          à partir de l'expérience de leurs auteurs, sans validation
          empirique de leurs niveaux ni de leurs critères.
    \item \textbf{La fiabilité de la mesure.} Deux évaluateurs n'attribuent
          pas toujours le même niveau au même processus, surtout en
          autoévaluation.
    \item \textbf{Le lien avec la performance.} Qu'un niveau de maturité plus
          élevé produise de meilleurs résultats est postulé plus souvent que
          démontré. Certaines études l'observent~\cite{simonsson2010maturity},
          d'autres non, et le lien dépend du contexte.
\end{enumerate}

\begin{boiteRemarque}{Faut-il renoncer aux modèles de maturité~?}
Non, mais il faut les utiliser pour ce qu'ils sont~: des outils de
\emph{diagnostic} et de \emph{dialogue}, pas des instruments de mesure
scientifique. COBIT~2019 en tient compte de deux manières, déjà vues dans
ce cours~: il fait dépendre la capacité \emph{cible} des facteurs de
conception, ce qui répond à la critique de la linéarité
(Chapitre~\ref{chap:cobit-conception}), et il rattache chaque niveau à des
activités précises, ce qui améliore la fiabilité de l'évaluation
(Chapitre~\ref{chap:cobit-capacite}). Pöppelbuß et Röglinger ont proposé des
principes de conception qui permettent de juger si un modèle de maturité
est utile~\cite{poppelbuss2011maturity}.
\end{boiteRemarque}

% ============================================================================
\section{Lire une Étude Empirique}
\label{sec:c20-lecture}
% ============================================================================

\subsection{Une grille de lecture}

\begin{tableaucours}{L{3.2cm}L{9.0cm}}{%
    \tableauHeader{Élément} &
    \tableauHeader{Questions à se poser}
}
Question et théorie & Quelle question précise~? Quelle théorie justifie
les hypothèses (agence, alignement, contingence)~? \\
Échantillon & Combien d'organisations, de quels pays, secteurs, tailles~?
Comment ont-elles été choisies~? \\
Mesures & Comment la gouvernance et la performance sont-elles mesurées~?
Par qui~? Les échelles sont-elles validées~? \\
Biais de méthode & Le même répondant décrit-il à la fois la gouvernance et
la performance~? \\
Causalité & L'étude est-elle transversale (une seule date) ou
longitudinale~? Peut-elle établir le sens de l'effet~? \\
Généralisation & Les résultats valent-ils au-delà de l'échantillon~? Pour
une organisation camerounaise~? \\
Statut & Article évalué par des pairs, communication, prépublication,
rapport de cabinet~? \\
\end{tableaucours}
\vspace{-0.8em}
\begin{center}{\small\itshape Tableau~20.3. Grille de lecture critique d'une
étude empirique}\end{center}

\subsection{Le statut des sources}

\begin{boxdef}
\termecle{Évaluation par les pairs}\index{évaluation par les pairs}
(\textit{peer review})~: examen d'un article par des chercheurs
indépendants avant sa publication dans une revue ou des actes de
conférence. Elle ne garantit pas que les résultats sont justes, mais
qu'ils ont été soumis à une critique méthodologique.

\medskip
\termecle{Prépublication}\index{prépublication} (\textit{preprint})~:
version d'un travail rendue publique avant son évaluation par les pairs.
Elle doit être citée comme telle.
\end{boxdef}

Les enquêtes des cabinets de conseil et des éditeurs de logiciels sont
utiles pour connaître les pratiques, mais elles poursuivent souvent un
objectif commercial et publient rarement leur méthode. Elles se citent
avec cette réserve.

\begin{boxscenario}{Exercice de lecture. Une étude fictive}
\textbf{Résumé.} « Une enquête menée auprès de 85~DSI d'entreprises
camerounaises montre que les entreprises qui ont adopté COBIT ont une
satisfaction client supérieure de 15~\% à celles qui ne l'ont pas adopté
($p < 0{,}05$). »

\medskip
\textbf{Lecture critique.} L'échantillon est-il représentatif (85~DSI
volontaires, sans doute les plus sensibles au sujet)~? La satisfaction
client est-elle mesurée auprès des clients, ou estimée par le DSI~? Les
entreprises qui adoptent COBIT sont-elles plus grandes, plus riches, mieux
organisées par ailleurs~? Une différence statistiquement significative
n'est pas une preuve d'effet. Conclusion~: résultat intéressant, à
présenter comme une association observée sur un échantillon de
convenance.
\end{boxscenario}

% ============================================================================
\section{Des Pistes de Recherche en Afrique Centrale}
\label{sec:c20-afrique}
% ============================================================================

\subsection{Un champ peu exploré}

La recherche publiée sur la gouvernance des SI porte majoritairement sur
des organisations des pays industrialisés. Les travaux consacrés à
l'Afrique centrale restent peu nombreux, alors que le contexte y présente
des caractéristiques propres~: concentration du pouvoir de décision,
tutelle publique, dépendance aux intégrateurs étrangers, régulation
récente, contraintes d'infrastructure (Chapitre~\ref{chap:gouv-entreprise}).
C'est une raison suffisante pour que les mémoires et les thèses s'en
emparent.

\subsection{Des questions possibles}

\begin{tableaucours}{L{5.6cm}L{3.2cm}L{3.4cm}}{%
    \tableauHeader{Question} &
    \tableauHeader{Méthode adaptée} &
    \tableauHeader{Terrain}
}
Comment les droits de décision informatiques se répartissent-ils dans
les établissements publics camerounais~? & Études de cas multiples & Trois
à cinq établissements \\
Les comités stratégiques des SI créés après une crise fonctionnent-ils
réellement~? & Étude longitudinale & Une banque, sur trois ans \\
Quels facteurs de conception COBIT pèsent le plus dans les organisations
de la zone CEMAC~? & Enquête par questionnaire & Banques et assurances \\
Comment adapter un modèle de capacité aux PME camerounaises~? &
Recherche orientée conception~\cite{hevner2004design} & PME de Douala et
Yaoundé \\
L'entrée en vigueur de la loi n°~2024/017 modifie-t-elle la gouvernance
des données~? & Études de cas avant et après & Entreprises de paiement
mobile \\
\end{tableaucours}
\vspace{-0.8em}
\begin{center}{\small\itshape Tableau~20.4. Exemples de questions de
recherche}\end{center}

\subsection{Les méthodes}

\begin{itemize}
    \item \textbf{L'étude de cas} convient aux questions « comment~? » et
          « pourquoi~? » dans des contextes peu connus~; elle exige un
          accès de confiance au terrain et une rigueur dans la collecte des
          preuves.
    \item \textbf{L'enquête par questionnaire} permet de mesurer des
          associations~; elle exige des échelles validées et un échantillon
          suffisant, souvent difficile à réunir.
    \item \textbf{La recherche orientée conception} construit un artefact
          (un modèle, une méthode, un outil) et l'évalue en situation~; elle
          convient bien aux praticiens qui préparent un mémoire
          professionnel~\cite{hevner2004design}.
\end{itemize}

\begin{boxprocedure}
\textbf{Un projet de mémoire recevable.} Le projet présenté au directeur
de mémoire comporte~:
\begin{itemize}
    \item une question de recherche précise, et sa justification au regard
          de la littérature~;
    \item le cadre théorique mobilisé~;
    \item la méthode, le terrain et l'accès obtenu (accord écrit de
          l'organisation)~;
    \item les règles de confidentialité et d'anonymisation des données
          recueillies~;
    \item le calendrier et les livrables~;
    \item une bibliographie d'au moins quinze références, dont une majorité
          évaluées par les pairs.
\end{itemize}
\end{boxprocedure}

\begin{boxmetier}{Auditeur SI}
\textbf{Faire de sa pratique un terrain de recherche}
\begin{enumerate}
    \item Les missions d'audit produisent des données rares sur la
          gouvernance réelle des organisations~: écarts, causes,
          capacités.
    \item Elles ne peuvent servir à la recherche qu'avec l'accord des
          organisations et une anonymisation rigoureuse, dans le respect de
          la confidentialité (Chapitre~\ref{chap:audit}).
    \item Agrégées sur plusieurs missions, elles permettent de répondre à
          des questions qu'aucune enquête par questionnaire n'atteint.
\end{enumerate}
\end{boxmetier}

% ============================================================================
\separateur
% ============================================================================

\begin{boxresume}
\textbf{Ce qu'il faut retenir de ce chapitre~:}
\begin{itemize}
    \item \textbf{Trois courants}~: les structures de décision, les
          mécanismes, l'effet sur la performance.
    \item \textbf{Des résultats convergents}~: la gouvernance compte, les
          mécanismes relationnels comptent autant que les structures, et
          l'effet passe par l'alignement.
    \item \textbf{Des questions ouvertes}~: le sens de la causalité, la
          fiabilité des mesures, la transposition des résultats, l'effet
          propre des référentiels.
    \item \textbf{Les modèles de maturité} sont des outils de diagnostic et
          de dialogue, pas des instruments scientifiques~: linéarité,
          fondement, fiabilité et lien avec la performance sont discutés.
    \item \textbf{Lire une étude}, c'est en examiner l'échantillon, les
          mesures, les biais, la causalité et le statut.
    \item \textbf{L'Afrique centrale} est un terrain peu exploré, qui
          appelle études de cas, enquêtes et recherche orientée conception.
\end{itemize}
\end{boxresume}

\bigskip

\begin{boxexo}[title={\ding{45}\quad Exercice~20.1. Association ou causalité~? \niveaubase}]
Pour chacune des affirmations suivantes, indiquez si elle décrit une
association ou un effet causal, et ce qu'il faudrait pour établir la
causalité.
\begin{enumerate}[(a)]
    \item « Les banques dotées d'un comité stratégique des SI ont moins
          d'incidents majeurs. »
    \item « Après la création de son comité, la BANQUE SANAGA a réduit ses
          incidents majeurs de moitié. »
    \item « Les entreprises certifiées ISO/IEC~27001 subissent moins de
          fuites de données. »
\end{enumerate}
\end{boxexo}

\begin{boxexo}[title={\ding{45}\quad Exercice~20.2. Lire un article \niveaumoyen}]
Choisissez l'un des articles cités dans ce chapitre et appliquez-lui la
grille du Tableau~20.3. Rédigez une note de lecture de deux pages~:
question, théorie, méthode, résultats, limites, intérêt pour une
organisation camerounaise.
\end{boxexo}

\begin{boxexo}[title={\ding{45}\quad Exercice~20.3. Un projet de mémoire \niveauexpert}]
Rédigez un projet de mémoire sur l'une des questions du Tableau~20.4, ou
sur une question de votre choix, en respectant l'exigence de contrôle de
la section~\ref{sec:c20-afrique}. Justifiez en particulier le choix de la
méthode et la manière d'obtenir l'accès au terrain.
\end{boxexo}

\bigskip

\begin{boxinfo}
\textbf{Transition.} Ce chapitre a présenté l'état de la recherche et les
questions qu'elle laisse ouvertes. Le Chapitre~\ref{chap:contributions}
présente des contributions et des pistes de recherche propres à
l'auteur et à son laboratoire.
\end{boxinfo}
